\selectlanguage{romanian}
\chapter*{Rezumat}
\addcontentsline{toc}{chapter}{Rezumat}

Lucrarea analizeaza proiectarea, implementarea si evaluarea platformei Trainee, o aplicatie full-stack dezvoltata pentru a intermedia relatia dintre clienti si antrenori personali. Solutia include o aplicatie mobila realizata cu Expo Router si React Native, un backend REST API construit cu Express si TypeScript, persistenta relationala in PostgreSQL si integrare cu servicii externe pentru plati, email si stocare media.

Contributia centrala este definirea unei arhitecturi modulare, capabile sa sustina fluxuri functionale complete: cautare de antrenori, programare sesiuni, interactiune cu sali de sport, raportare incidente si monetizare prin subscrieri. In paralel, lucrarea trateaza explicit dimensiunea de securitate: validare stricta a datelor de intrare, rate limiting pe endpoint-uri publice, control de acces pe roluri si management riguros al secretelor prin variabile de mediu.

Rezultatele obtinute confirma functionarea coerenta a fluxurilor end-to-end pentru rolurile de client, antrenor si administrator. In forma actuala, platforma ofera un fundament tehnic stabil si extensibil, adecvat pentru etape ulterioare de observabilitate, optimizare si scalare.

\textbf{Cuvinte cheie:} aplicatie full-stack, React Native, Express, PostgreSQL, securitate API, scheduling, marketplace fitness.

\clearpage

\selectlanguage{english}
\chapter*{Abstract}
\addcontentsline{toc}{chapter}{Abstract}

This thesis examines the design, implementation, and evaluation of Trainee, a full-stack platform built to connect clients with personal trainers in a unified digital product. The implemented solution combines a mobile application based on Expo Router and React Native, an Express and TypeScript backend API, a PostgreSQL persistence layer, and third-party integrations for payments, email communication, and media storage.

The core contribution is a modular architecture that supports complete functional flows, including trainer discovery, schedule management, gym-related interactions, incident reporting, and subscription-based billing. The work also highlights practical security choices: strict request validation, endpoint-specific rate limiting, role-based access control, and secure secret handling through environment configuration.

The evaluation results indicate that the platform supports coherent end-to-end flows for client, trainer, and administrator roles, while preserving a solid and extensible technical base for future observability, monitoring, and scaling improvements.

\textbf{Keywords:} full-stack application, React Native, Express, PostgreSQL, API security, scheduling, fitness marketplace.

\clearpage
\selectlanguage{romanian}
